\documentclass[12pt,a4paper]{article}

\usepackage[margin=1in]{geometry}
\usepackage[T1]{fontenc}
\usepackage[utf8]{inputenc}
\usepackage{lmodern}
\usepackage{enumitem}
\usepackage{hyperref}

\title{Process \& Collaboration (10\%)\\Jobra Hospital Management System}
\author{}
\date{\today}

\begin{document}
\maketitle

\section{Agile methodology adherence}
\begin{itemize}[leftmargin=*, itemsep=0.4em]
  \item \textbf{Incremental delivery}: The repository history reflects milestone-style progress (initial setup, bulk upload, near-completion, documentation addition, final report PDF addition), which aligns with iterative development.
  \item \textbf{Backlog-style improvements captured}: The README includes a ``Future Improvements'' section, which functions like a lightweight backlog for planned enhancements.
\end{itemize}

\section{Team collaboration effectiveness}
\begin{itemize}[leftmargin=*, itemsep=0.4em]
  \item \textbf{Clear role boundaries}: The system is partitioned by role (admin/doctor/patient) and by module (appointments, reports, complaints), which makes it easier for multiple team members to work in parallel on different areas.
  \item \textbf{Shared conventions}: Common patterns (route prefixing under \texttt{/api}, shared auth approach, centralized API client in frontend) improve team coordination and reduce merge conflicts.
\end{itemize}

\section{Version control usage}
\begin{itemize}[leftmargin=*, itemsep=0.4em]
  \item \textbf{Repository tracked in Git}: Commits exist for major checkpoints, including adding the README and final documentation artifacts.
  \item \textbf{Improvement note}: More granular commits (smaller, feature-scoped messages) and use of branches/pull requests would strengthen traceability and review quality for a larger team.
\end{itemize}

\section{Meeting participation}
\begin{itemize}[leftmargin=*, itemsep=0.4em]
  \item \textbf{Suggested evidence to include in viva/report}: Meeting minutes, sprint planning notes, task assignments, and short weekly progress summaries can be added as appendices or links to demonstrate participation and coordination.
  \item \textbf{Practical workflow}: A regular cycle of planning (tasks), implementation (coding), review (demo/testing), and retrospective (improvements list) matches the structure of this project and is easy to present academically.
\end{itemize}

\end{document}

